\subsection{Part A – Concept Check}

\indent 

\noindent\textbf{Why is information considered the “fuel” of IT projects?}

Information guides decision-making at every stage of an IT project, from planning to implementation. Without accurate information, project managers cannot define scope, allocate resources, or assess risks. Example: Knowing system requirements, user expectations, and technology constraints enables efficient project execution.

\bigskip

\noindent\textbf{Explain the DIKW hierarchy using an IT-related example (e.g., social media data $\xrightarrow{}$ analytics $\xrightarrow{}$ business insights).}


Data: Raw user activity logs from a social media platform. 

Information: Processed data showing the number of daily active users and time spent per session.

Knowledge: Analyzing trends to identify peak usage times and user engagement patterns. 

Wisdom: Using insights to make strategic decisions, such as optimizing app notifications or server load balancing to improve user experience.


\bigskip

\noindent\textbf{Why are statistics important when conducting IT research (e.g., system load testing, customer survey results)?}


Statistics summarize and interpret large volumes of data to identify patterns, trends, and anomalies. Descriptive statistics (mean, median, variance) help understand system performance or user behaviour. Inferential statistics enable predictions and testing hypotheses, e.g., evaluating if a new feature improves user engagement.

\subsection{Part B – Case Study continuation}

\indent 

Case Study: MediCare Solutions EHR Project A mid-sized healthcare company, MediCare Solutions, wants to build an Electronic Health Record (EHR) system with:

\begin{itemize}
    \item Patient records storage (structured + unstructured)
    \item Appointment scheduling
    \item Billing system integration
    \item A mobile app for patients
\end{itemize}

The project manager must estimate the size, effort, resources, duration, and cost. They must also ensure information validity and use the right research methods.

\bigskip

\noindent\textbf{Question:}

\bigskip

\noindent\textbf{Apply the DIKW model to healthcare data in this project.}

Data: Raw patient records, appointment logs, billing entries.

Information: Processed reports showing patient visits per department or payment status.

Knowledge: Patterns like peak appointment hours, most common procedures, or billing delays.

Wisdom: Using these insights to optimize scheduling, resource allocation, and improve patient care.

\bigskip

\noindent\textbf{Identify 2 primary research and 2 secondary research sources relevant here.}

Primary research: Surveys/interviews with hospital staff to understand workflow needs. Observations of patient interactions and appointment handling processes.

Secondary research: Industry reports on EHR system adoption and standards (e.g., HIMSS reports). Academic papers or case studies on healthcare IT implementations.


\bigskip

\noindent\textbf{Discuss why credible and referenced information is especially critical for healthcare IT projects.}


Ensures compliance with regulations (HIPAA, GDPR, local healthcare laws). Reduces risk of errors that could compromise patient safety or data security. Provides reliable input for system design, testing, and decision-making.

\bigskip

\noindent\textbf{Using the 5 steps of estimation, make a rough plan}


\step Estimate project size (modules/features).
\step Effort required (person-months).
\step Resources (number of developers, testers, etc.).
\step Duration (months).
\step Cost (assume \$8,000/month per developer, \$6,000/month per tester).


\bigskip

\noindent\textbf{Which estimation approach (expert judgment, analogy, function points, empirical models) seems best for this case, and why?}
